\documentclass[preview]{standalone}

\usepackage{amsmath}
\usepackage{amssymb}
\usepackage{bettelini}
\usepackage{stellar}
\usepackage{definitions}

\begin{document}

\id{measuretheory-lebesgue-regularity}
\genpage

\section{Borel and Lebesgue Measurability}

\begin{snippettheorem}{borel-restriction-lebesgue-measure-not-complete}{The Borel restriction is not complete}
    The restriction of Lebesgue measure to the Borel \(\sigma\)-algebra,
    \(\restr{\mu}{\mathcal B(\realnumbers)}\), is not complete.
\end{snippettheorem}

\begin{snippetproof}{borel-restriction-lebesgue-measure-not-complete-proof}{borel-restriction-lebesgue-measure-not-complete}{The Borel restriction is not complete}
    The Cantor set is a Borel set of measure zero and has the cardinality of
    the continuum. Its power set therefore has cardinality strictly larger
    than that of \(\mathcal B(\realnumbers)\). Consequently, some subset of
    the Cantor set is not Borel. Such a subset is nevertheless Lebesgue
    measurable because Lebesgue measure is complete.
\end{snippetproof}

In fact, the Lebesgue \(\sigma\)-algebra is precisely the completion of the
Borel \(\sigma\)-algebra with respect to Lebesgue measure. The following
regularity result makes this explicit.

\begin{snippetproposition}{lebesgue-measurable-set-regularity}{Regularity of Lebesgue measurable sets}
    Let \(E\in\mathcal L\). Then:
    \begin{enumerate}
        \item there is a \(G_\delta\) set \(G\supseteq E\) such that
            \(\mu(G\setminus E)=0\);
        \item there is an \(F_\sigma\) set \(F\subseteq E\) such that
            \(\mu(E\setminus F)=0\);
        \item if \(\mu(E)<+\infty\), then for every \(\varepsilon>0\) there
            are bounded intervals \(I_1,\ldots,I_N\) such that
            \[
                \mu\!\left(E\mathbin{\triangle}\bigcup_{j=1}^N I_j\right)
                <\varepsilon.
            \]
    \end{enumerate}
\end{snippetproposition}

Thus, up to a null set, every Lebesgue measurable set lies between Borel sets.
Conversely, if \(E\subseteq G\in G_\delta\) and
\(\mu^*(G\setminus E)=0\), then \(E\in\mathcal L\).

\begin{snippetproof}{lebesgue-measurable-set-regularity-proof}{lebesgue-measurable-set-regularity}{Regularity of Lebesgue measurable sets}
    First suppose that \(\mu(E)<+\infty\). By the definition of outer measure,
    for every \(\varepsilon>0\) there are bounded open intervals \(I_n\) such
    that
    \[
        E\subseteq A\triangleq\bigcup_{n=1}^{\infty}I_n,
        \qquad
        \sum_{n=1}^{\infty}\ell(I_n)<\mu(E)+\varepsilon.
    \]
    Since \(E\) is measurable and \(E\subseteq A\),
    \[
        \mu(A)=\mu(E)+\mu(A\setminus E),
    \]
    so \(\mu(A\setminus E)<\varepsilon\).

    For every \(n\), choose such an open set \(A_n\supseteq E\) with
    \(\mu(A_n\setminus E)<1/n\), and put
    \[
        G=\bigcap_{n=1}^{\infty}A_n.
    \]
    Then \(G\in G_\delta\), \(G\supseteq E\), and
    \[
        \mu(G\setminus E)\leq\mu(A_n\setminus E)<\frac1n
    \]
    for every \(n\); hence \(\mu(G\setminus E)=0\).

    If \(\mu(E)=+\infty\), split
    \[
        E=\bigcup_{k\in\integers}\bigl(E\cap[k,k+1]\bigr).
    \]
    For a fixed error \(\varepsilon\), approximate the \(k\)-th bounded piece
    from above by an open set with error at most
    \(\varepsilon/2^{|k|+2}\), then take the union of these open sets. Repeating
    this construction with error \(1/n\) and finally intersecting the resulting
    open sets gives the required \(G_\delta\) set.

    Apply the first part to \(E^c\). If \(G\supseteq E^c\) is \(G_\delta\) and
    \(\mu(G\setminus E^c)=0\), then \(F=G^c\) is \(F_\sigma\), \(F\subseteq E\),
    and
    \[
        E\setminus F=G\setminus E^c,
    \]
    so \(\mu(E\setminus F)=0\).

    Finally suppose \(\mu(E)<+\infty\). Choose an open
    \(A=\bigcup_{n\geq1}I_n\supseteq E\) with
    \(\mu(A\setminus E)<\varepsilon/2\). Since
    \(\sum_n\ell(I_n)<+\infty\), choose \(N\) such that
    \(\sum_{n>N}\ell(I_n)<\varepsilon/2\). For
    \(U_N=\bigcup_{n=1}^NI_n\),
    \[
        E\mathbin{\triangle}U_N
        \subseteq(A\setminus E)\cup\bigcup_{n>N}I_n,
    \]
    and therefore \(\mu(E\mathbin{\triangle}U_N)<\varepsilon\).
\end{snippetproof}

The last assertion says that all the pathology of a finite-measure Lebesgue
set can be confined to a set of arbitrarily small measure. For instance, the
unbounded region under a Gaussian curve has finite area, so it can be
approximated by finitely many rectangles after discarding tails of arbitrarily
small area. This fails for an infinite-area strip.

\begin{snippetexample}{large-compact-subset-of-irrationals}{A large compact subset of the irrationals}
    Let
    \[
        E=[0,1]\setminus\rationalnumbers
    \]
    and enumerate the rationals as \(\{q_n\}_{n\geq1}\). For
    \(\varepsilon>0\), consider
    \[
        U=\bigcup_{n=1}^{\infty}
        \left(q_n-\frac{\varepsilon}{2^{n+1}},
        q_n+\frac{\varepsilon}{2^{n+1}}\right).
    \]
    This is an open set containing all rationals and
    \(\mu(U)\leq\varepsilon\). Hence
    \[
        K=[0,1]\setminus U
    \]
    is compact, contained in \(E\), and has measure at least
    \(1-\varepsilon\).
\end{snippetexample}

\end{document}
